\documentclass[conference,10pt]{IEEEtran}
\IEEEoverridecommandlockouts
\usepackage[utf8]{inputenc}
\usepackage[T1]{fontenc}
\usepackage{amsmath,amssymb}
\usepackage{booktabs}
\usepackage{tabularx}
\usepackage{array}
\usepackage{enumitem}
\usepackage[hidelinks]{hyperref}
\usepackage{microtype}
\usepackage{placeins}

\newcolumntype{Y}{>{\raggedright\arraybackslash}X}
\newcolumntype{L}[1]{>{\raggedright\arraybackslash}p{#1}}
\setlist{nosep,leftmargin=*}
\renewcommand{\arraystretch}{1.08}
\setlength{\tabcolsep}{3pt}

\begin{document}

\title{Bump Once, Remember Wisely: Calibrated Collision Detection and Forgetting Memory for Aerial Robots Among Invisible Obstacles}

\author{\IEEEauthorblockN{Egemen Çevik (150220012), Cumali Sarı (150230038), Ahmet Enes Çiğdem (150220079)}
\IEEEauthorblockA{YZV406E Robotics -- Term Project Proposal (Round 1), Istanbul Technical University, 10.11.2026}}

\maketitle

%==================================================================
\section{Problem Description and Primary Reference Paper}

\textbf{Problem.} Autonomous micro aerial vehicles (MAVs) navigate cluttered indoor spaces with LiDAR, assuming that every obstacle is seen. Two common cases break this assumption: (1)~\emph{invisible obstacles}---glass walls, doors and windows, transparent curtains and thin wires let laser beams pass through, so the map shows free space and the planner flies into them; and (2)~\emph{obstacles that enter unseen}, e.g., from above or below the LiDAR's limited vertical field of view. In both cases a collision cannot be avoided from perception alone. The MAV must therefore (i)~\emph{detect} that it collided, (ii)~\emph{react} safely before it crashes, and (iii)~\emph{remember} the obstacle so that it does not hit it again. Remembering creates a new difficulty: if the obstacle was a person, an object that is later moved, or a door that later opens, a permanent memory turns it into a \emph{phantom wall}. A \emph{false} detection does the same and can make the goal unreachable. The project addresses the suggested topic \emph{Robot Failure Detection and Recovery}. \textbf{Our problem} is to detect MAV collisions with a \emph{known} false-alarm rate, and to remember them selectively: keep invisible static obstacles, forget transient ones and false alarms.

\textbf{Primary reference paper.} R.~Wang \emph{et al.}, ``Air Bumper: A Collision Detection and Reaction Framework for Autonomous MAV Navigation,'' \emph{IEEE ICRA 2024}, pp.~15735--15741~\cite{airbumper}. It is an original research paper from an approved venue (ICRA) in the required period. A caged MAV with a 3D LiDAR uses only its IMU to: \emph{detect} a collision when any body-axis acceleration exceeds a fixed threshold ($a^*=20$\,m/s$^2$, sliding window $N=10$); \emph{estimate} its direction and contact point on a spherical cage; \emph{react} immediately with a position setpoint away from the hit, blended with the gradient of the Euclidean distance transform (EDT) of known obstacles; and \emph{remember} it by inserting the contact patch into the volumetric map as a ``collision sensor'' whose marks are \emph{never cleared} by the LiDAR. Recovery from glass doors is shown in Gazebo, and from glass and stick hits on a real MAV (detection $\approx$0.18\,s, recovery $<$1\,s, 10 trials).

\textbf{Why this paper.} It is an early framework that combines collision detection, reaction \emph{and} mapping in a standard LiDAR navigation stack using only the IMU. It turns failures into map knowledge (``failure as a sensor''), which is directly a \emph{failure detection and recovery} problem. It is open source~\cite{airbumpercode} and validated in simulation and on hardware. Since its preprint (July 2023) it has been cited by follow-up work on contact-aware flight, e.g., whisker-based tactile navigation in \emph{Nature Communications}~\cite{whisker} and a survey of LiDAR-based aerial vehicles in \emph{IEEE/ASME T-Mech}~\cite{lidarsurvey} (7 citing papers on Semantic Scholar, Nov.~2026). Its design choices are simple and explicit (fixed threshold, permanent marks), so their limitations can be analyzed and improved within a term project. Related work handles failures differently: \cite{gp} predicts planner failures \emph{before} they happen from LiDAR geometry, and \cite{ifar} adapts edge costs online from contact feedback; both assume the LiDAR sees every obstacle, which is exactly what fails with glass.

\textbf{Significance.} Indoor MAVs for inspection, warehouse inventory and search and rescue operate around glass facades, windows and partitions, and among people and moving objects. A collision detector whose false-alarm rate is unknown, and a map that never forgets, both cause mission failures: false reactions interrupt flight, and phantom walls block routes. We keep the paper's platform (a caged MAV) and its detect--react--remember pipeline, and improve the two weakest parts: \emph{how} a collision is declared and \emph{how long} it is remembered.

%==================================================================
\section{Environment, Benchmark, and Dataset}

\begin{table}[t]
\caption{Software and simulation environment (all free).}
\label{tab:env}
\footnotesize
\begin{tabularx}{\columnwidth}{@{}L{1.45cm}YY@{}}
\toprule
\textbf{Component} & \textbf{Choice} & \textbf{Justification}\\
\midrule
Middleware & ROS~2 Jazzy~\cite{ros2} in Docker & LTS, free; identical setup on every laptop\\
Autopilot & PX4 SITL~\cite{px4} via uXRCE-DDS (\texttt{px4\_msgs}), offboard position/velocity setpoints & Same autopilot as~\cite{airbumper}; reaction setpoints go directly to the low-level controller\\
Simulator & Gazebo Harmonic~\cite{gazebo} (PX4 \texttt{gz} integration), headless batch runs & Same simulator family as~\cite{airbumper}; free ground truth; \emph{WindEffects} for gusts\\
Vehicle & PX4 x500 quadrotor with a spherical collision cage (radius $l=0.25$\,m), 3D LiDAR configured like a Livox Mid-360 (360$^\circ\times$59$^\circ$), IMU (250\,Hz, noise + bias), cage contact sensor & Mirrors the paper's caged MAV; contact sensor gives ground-truth labels and is \emph{not} used by the detector\\
Map \& planner & Local voxel occupancy map (0.1\,m) with EDT~\cite{edt}, A$^*$ path search, velocity-limited setpoint follower & Lightweight Python replacement for the paper's GPU mapper and GTO-MPC; exposes cost and speed limits per voxel\\
Invisible obstacles & Models with collision but no visual geometry & GPU LiDAR renders visuals only, so beams pass through while physics blocks: glass as in~\cite{airbumper}\\
Moving / transient & Joint/trajectory-driven panels and poles, doors opening at random times, objects removed after contact, a ``stick'' hitting from outside the LiDAR FOV & Creates ``remembered but gone'' cases and the paper's stick test\\
\bottomrule
\end{tabularx}
\end{table}

Table~\ref{tab:env} summarizes the environment. Because no public benchmark contains invisible and transient obstacles for MAVs, we build a procedurally generated, seeded \textbf{scenario suite}, including the paper's own layout:
\begin{itemize}
\item \textbf{S1 -- Static glass} (60 worlds): indoor rooms and corridors with 1--4 glass walls, glass doors (re-creating the three-door world of~\cite{airbumper}) and windows at random positions and heights. Tests detection and memory on invisible static obstacles.
\item \textbf{S2 -- Transient} (45): (a)~a glass door that opens 20--120\,s after the first contact; (b)~a visible panel blocking the shortest route, removed after contact; (c)~a pole hitting the cage from below or above the LiDAR FOV (the paper's stick test). Tests the phantom-wall problem.
\item \textbf{S3 -- False-alarm stress} (30): \emph{no collisions}; wind gusts (WindEffects), aggressive stop-and-go and sharp turns, high-speed flight, doubled IMU noise. Measures false alarms per hour.
\end{itemize}
Worlds are split 50/20/30\% into train/calibration/test per set; gust and maneuver profiles unseen during calibration form an out-of-distribution test. Speeds $v_{\max}\in\{0.5,1.0,1.5\}$\,m/s. We plan $\sim$10\,h of collision-free flight for calibration (Sec.~\ref{sec:method}), a \emph{separate} collision-free test set of $\ge$50\,h to verify false-alarm rates (headless runs are cheap), and $\sim$2{,}000 labeled collisions, which fits into overnight runs with 2--3 parallel containers.

%==================================================================
\section{Challenges and Limitations}

Table~\ref{tab:lim} lists the limitations of~\cite{airbumper} that motivate our project, each with evidence from the paper.

\begin{table}[t]
\caption{Limitations of Air Bumper~\cite{airbumper}.}
\label{tab:lim}
\footnotesize
\begin{tabularx}{\columnwidth}{@{}L{0.45cm}L{2.55cm}Y@{}}
\toprule
\# & \textbf{Limitation (evidence)} & \textbf{Why it matters}\\
\midrule
L1 & Fixed, hand-tuned threshold on raw acceleration, $a^*=20$\,m/s$^2$ (Sec.~IV-A, ``experimentally determined''); no false-alarm or miss rates reported & The MAV's \emph{own} aggressive maneuvers and wind gusts also cause spikes, while soft or low-speed hits stay below $a^*$. The false-alarm rate is unknown and changes with platform, speed and wind\\
L2 & Per-axis test $|a_x|>a^*$ \emph{or} $|a_y|>a^*$ (Sec.~IV-A) & A diagonal hit splits between axes: 25\,m/s$^2$ at 45$^\circ$ gives 17.7\,m/s$^2$ per axis and is \emph{missed}\\
L3 & Direction from raw acceleration, $\mathbf{a}_c=-\mathbf{a}$ (Eqs.~1--4) & Commanded acceleration biases the estimated direction, and thus the recovery setpoint and the mapped patch, especially while braking or turning\\
L4 & Permanent, top-priority marks: ``the robot remembers all obstacles it ever collides with'' (Sec.~V-C) & Right for static glass; wrong for people, moved objects, opened doors and \emph{false detections}, which become phantom walls. Detection errors (L1) are made \emph{permanent} by the map\\
L5 & LiDAR evidence at the contact is ignored & Whether the LiDAR \emph{saw} the obstacle separates visible (already mapped, maybe moving) from invisible (must be remembered)\\
L6 & Small, qualitative evaluation: 1 simulated world, 10 real trials & No TPR, false alarms/hour, latency or direction error; no baselines or statistics\\
\bottomrule
\end{tabularx}
\end{table}

These issues matter because a failure detector is only deployable when its false-alarm and miss rates are known, and because a map that never forgets turns every detection error and every moved object into a lasting navigation failure. The open questions are: \emph{Can MAV collision detection have a guaranteed false-alarm rate instead of a hand-tuned threshold? Can collision memory be selective, keeping invisible static obstacles while forgetting transient ones and its own mistakes?}

%==================================================================
\section{Proposed Method and Novel Contributions}
\label{sec:method}

\textbf{Baseline reimplementation (B1, not a contribution).} We first port Air Bumper to ROS~2 and PX4 SITL: the per-axis detector ($a^*$, $N$), the sphere contact model $\mathbf{c}_c=\hat{\mathbf{C}}l$, the recovery setpoint $\mathbf{p}_r=(w\,\mathbf{G}/\|\mathbf{G}\|-(1-w)\hat{\mathbf{C}})R_d$ with $R_d=2l$ sent directly as a PX4 offboard position setpoint, and permanent, lethal collision voxels in the map.

\subsection{C1: Model-aware, calibrated collision detection (L1--L3)}
\textbf{(a) Residual signal.} Instead of raw acceleration we use the \emph{unexpected} acceleration
\begin{equation}
\mathbf{r}(t)=\mathbf{a}^{\text{IMU}}(t)-\hat{\mathbf{a}}(t),
\end{equation}
where $\hat{\mathbf{a}}$ is the acceleration explained by the MAV's own commands: the commanded velocity/acceleration setpoints are passed through a closed-loop response model of the PX4 controller identified on collision-free flights, and rotated into the body frame with gravity removed. Own braking and turning are thus removed (L1, L3). Model-based residuals are the standard tool for collision detection on manipulators~\cite{haddadin}; we bring this idea to IMU-based detection on MAVs.

\textbf{(b) Features and score.} Over a sliding window $W=0.1$\,s we compute the peak residual \emph{norm} $\|\mathbf{r}\|$ in 3D (removes the per-axis issue, L2), the peak jerk, the angular-rate residual, and the setpoint-tracking error (pushing against an obstacle). A small, interpretable logistic regression maps them to a score $s\in[0,1]$; as an ablation, $s=\|\mathbf{r}\|$.

\textbf{(c) Calibrated threshold.} We set the threshold by split conformal calibration~\cite{vovk,angelopoulos} on $n$ \emph{collision-free} windows from the calibration worlds:
\begin{equation}
\hat\tau = s_{(\lceil (n+1)(1-\alpha)\rceil)},
\end{equation}
the $\lceil (n+1)(1-\alpha)\rceil$-th smallest score. Under exchangeability a new collision-free window triggers with probability $\le\alpha$, so the false-alarm rate is fixed in advance: $\approx\alpha\cdot3600/W$ per hour for non-overlapping windows. For a target of 1 false alarm/h with $W=0.1$\,s, $\alpha\approx2.8\cdot10^{-5}$ and $n\ge1/\alpha\approx36{,}000$ windows ($\approx$1\,h of flight; we collect $\sim$10\,h). We target $\{1,5,20\}$ false alarms/h. Since windows are temporally correlated, we verify the rate empirically on the $\ge$50\,h test set with exact Poisson (Clopper--Pearson-type) 95\% intervals; at 1/h this gives $\sim$50 expected events.

\textbf{(d) Residual-based estimation.} The direction $\hat{\mathbf{C}}$ (azimuth and polar angle, as in~\cite{airbumper}) is computed from $-\mathbf{r}$ instead of $-\mathbf{a}$, so the recovery setpoint and the mapped contact patch are not biased by the MAV's own acceleration (L3).

\subsection{C2: Evidence-aware, forgetting collision memory (L4, L5)}
\textbf{(a) LiDAR evidence.} At detection we inspect the LiDAR points within a $\pm10^\circ$ cone around the contact direction and label the contact \emph{visible} (a return within contact distance $+0.15$\,m: already mapped, possibly a person or a moving object), \emph{invisible} (beams pass beyond the contact: the glass case), or \emph{unseen} (direction outside the vertical FOV, e.g., the stick from below).

\textbf{(b) Decaying belief.} Each mark $k$ stores its contact patch, hit count $h_k$ and last-hit time $t_k$, with belief
\begin{equation}
b_k(t)=\exp\!\big(-(t-t_k)/\tau_{\text{class}}\big),
\end{equation}
$\tau_{\text{invisible}}=600$\,s, $\tau_{\text{unseen}}=60$\,s, $\tau_{\text{visible}}=10$\,s (enough to prevent an immediate re-hit; the LiDAR covers the rest). All values are tuned on calibration worlds and studied for sensitivity. A new hit within 0.3\,m merges with the mark ($h_k{+}1$, $t_k$ reset): a real glass wall is hit again and stays remembered, while a transient obstacle or a false alarm fades.

\textbf{(c) Graded costs and speed limits, not walls.} A collision voxel gets planning cost $c=c_{\max}b_k$, and is \emph{lethal only if} $h_k\ge2$ (confirmed); otherwise it is high but passable, so a single unconfirmed mark can never make the goal unreachable. Near decayed or unconfirmed marks the setpoint speed is limited to 0.3\,m/s, so re-testing a remembered spot causes at most a gentle cage contact that C1 detects quickly.

\textbf{Implementation.} All nodes are Python (\texttt{rclpy}): \texttt{imu\_residual}, \texttt{collision\_detector}, \texttt{reaction}, \texttt{collision\_memory}, \texttt{voxel\_planner}. Memory only changes voxel costs and speed limits, so it can be added to any voxel-based planner. \textbf{Expected improvements over~\cite{airbumper}:} (1)~a false-alarm rate fixed before deployment and verified on test worlds, robust to wind and aggressive maneuvers; (2)~detection of diagonal and low-speed contacts; (3)~smaller direction errors and better-placed recovery setpoints; (4)~no phantom walls with transient obstacles or false alarms, at the cost of a small, measured increase in repeated collisions on static glass (a mark becomes lethal only after a second hit). Whether this trade-off is favorable is tested in RQ3.

%==================================================================
\section{Research Questions and Evaluation Plan}

\textbf{RQ1 (Detection):} Does C1 achieve a higher true-positive rate than Air Bumper's fixed per-axis threshold \emph{at the same false-alarm rate}, across speeds, maneuvers and wind, and does the calibrated rate hold under unseen gust profiles? \textbf{RQ2 (Estimation):} How much does residual-based estimation reduce the direction error and improve the recovery distance? \textbf{RQ3 (Memory):} Does C2 avoid phantom-wall failures with transient obstacles and false alarms, and how many extra repeated collisions with static invisible obstacles does this cost compared with permanent memory? \textbf{RQ4 (System):} Does the full system improve mission success, crashes and time-to-goal over the baselines?

\begin{table}[t]
\caption{Baselines and ablations.}
\label{tab:base}
\footnotesize
\begin{tabularx}{\columnwidth}{@{}L{0.75cm}YL{1.55cm}L{1.9cm}@{}}
\toprule
\textbf{ID} & \textbf{Detection} & \textbf{Memory} & \textbf{Represents}\\
\midrule
B0 & none (LiDAR-only navigation) & none & Standard MAV stack (crashes, as in~\cite{airbumper} Fig.~4b)\\
B1 & Air Bumper~\cite{airbumper}: $a^*=20$\,m/s$^2$ and $a^*$ re-tuned on calibration worlds & permanent, lethal & Primary paper\\
B2 & Oracle (cage contact sensor) & permanent & Perfect detection; isolates memory\\
B3 & C1 & permanent & C1 alone\\
B4 & Air Bumper & C2 & C2 alone\\
Ours & C1 & C2 & Full method\\
\bottomrule
\end{tabularx}
\end{table}

\textbf{Baselines} are listed in Table~\ref{tab:base}. \textbf{Metrics.} \emph{Detection} (per collision event, ground truth from the cage contact sensor): TPR, false alarms per hour of collision-free flight (with exact Poisson 95\% intervals), empirical vs.\ target false-alarm rate, detection latency, ROC-AUC. \emph{Estimation:} direction error (deg), post-recovery distance to the obstacle (m). \emph{Memory:} phantom volume (m$^3$ of marked cost with no real obstacle, integrated over time), repeated collisions (contact within 0.5\,m of an earlier one), blocked-goal events. \emph{System:} mission success, crashes (attitude $>60^\circ$ or uncommanded altitude loss), collisions per run, time-to-goal, node run time ($<$5\,ms per cycle).

\begin{table}[t]
\caption{Planned experiments.}
\label{tab:exp}
\footnotesize
\begin{tabularx}{\columnwidth}{@{}L{0.5cm}YL{1.6cm}L{0.75cm}@{}}
\toprule
\# & \textbf{Experiment} & \textbf{Data} & \textbf{RQ}\\
\midrule
E1 & Offline detection: ROC, TPR at matched false-alarm rate; features vs.\ $\|\mathbf{r}\|$ vs.\ raw acceleration; per speed and wind level & S1--S3 test logs & RQ1\\
E2 & Calibration validity: target $\{1,5,20\}$/h vs.\ measured; in-distribution vs.\ unseen gusts/maneuvers & $\ge$50\,h free test set & RQ1\\
E3 & Estimation: raw vs.\ residual direction; recovery distance & S1, S2 & RQ2\\
E4 & Memory (closed loop, oracle detector): permanent vs.\ decay-only vs.\ C2 & S1, S2 & RQ3\\
E5 & Full comparison B0--B4 vs.\ Ours $\times$ 3 speeds $\times$ 3 seeds & S1--S3 test & RQ4\\
E6 & Ablations/sensitivity: $\tau$, lethal-after-$h$, speed limit, LiDAR evidence, $W$ & Calibration worlds & RQ3, RQ4\\
\bottomrule
\end{tabularx}
\end{table}

\textbf{Experiments} are listed in Table~\ref{tab:exp}. \textbf{Compute budget.} With $\sim$41 test worlds, E5 has $6\times3\times3\times41\approx2{,}200$ episodes of $\sim$60\,s simulated flight ($\approx$37\,h); with PX4 SITL at a real-time factor $\approx$1 and 3 parallel containers this is $\approx$12\,h of wall-clock time. Fallback: one seed or two speeds for the full grid, keeping all seeds for B1 vs.\ Ours. \textbf{Statistics.} The world is the unit of analysis: 95\% bootstrap confidence intervals over worlds, McNemar tests for success/failure and Wilcoxon signed-rank tests for continuous metrics, with Holm correction. All hyperparameters, including the B1 threshold, are tuned only on calibration worlds. \textbf{Threats to validity.} Simulated contacts and IMU are idealized, so we add noise/bias and report results per noise level; the conformal guarantee is approximate under temporal correlation, so we report the empirical rate; sim-to-real transfer is not tested. Implementation risks are checked in week~1: whether collision-only models are invisible to the Gazebo \texttt{gpu\_lidar} sensor (fallback: a LiDAR visibility mask), and whether the caged x500 survives contacts in SITL without flipping (fallback: tuned contact stiffness and a lower speed range).

%==================================================================
\section{Team Roles and Individual Contributions}

\begin{table}[!h]
\caption{Team roles.}
\label{tab:team}
\footnotesize
\begin{tabularx}{\columnwidth}{@{}L{1.35cm}L{1.6cm}Y@{}}
\toprule
\textbf{Member} & \textbf{Role} & \textbf{Responsibilities and deliverables}\\
\midrule
Egemen Çevik & Simulation, scenarios \& experiments & Docker/ROS~2/PX4 SITL/Gazebo bring-up; caged x500 with LiDAR, IMU, contact sensor; invisible-obstacle models; S1--S3 generators; voxel map and planner; episode runner, ground-truth labels; batch runs of E5\\
Cumali Sarı & Detection \& estimation (C1) & Air Bumper detector for B1; controller-response model, residual node, features, logistic regression, conformal calibration, residual-based direction estimation; E1--E3; statistics\\
Ahmet Enes Çiğdem & Reaction \& memory (C2) & Air Bumper recovery control and permanent memory for B1; collision-memory node, LiDAR evidence check, graded costs and speed limits; B3/B4 configs; E4, E6; figures\\
\bottomrule
\end{tabularx}
\end{table}

Roles are given in Table~\ref{tab:team}. Each member writes and presents the part for their own module; contributions are documented by per-member git history and a weekly \texttt{CONTRIBUTIONS.md}. \textbf{Timeline:} week~1 (11--16 Nov) PX4 SITL + ROS~2 running and verification of the invisible-obstacle model and IMU collision spikes; week~2 S1--S3, voxel planner and B1; week~3 data collection and controller-response model; week~4 C1, E1--E3; week~5 C2, E4; week~6 E5, E6, statistics; week~7 report and slides; presentation on 30.12.2026. \textbf{Scope:} minimum = B1 + C1 + E1--E3; target = + C2, E4--E5; stretch = a real-world-like LiDAR model with partial glass returns, and a low-speed contact study.

%==================================================================
\FloatBarrier
\begin{thebibliography}{99}
\footnotesize
\bibitem{airbumper} R.~Wang, Z.~Guo, Y.~Chen, X.~Wang, and B.~M. Chen, ``Air Bumper: A collision detection and reaction framework for autonomous MAV navigation,'' in \emph{Proc. IEEE Int. Conf. Robot. Autom. (ICRA)}, 2024, pp. 15735--15741, doi: 10.1109/ICRA57147.2024.10611410.
\bibitem{airbumpercode} R.~Wang \emph{et al.}, ``Air Bumper,'' GitHub repository, 2024. [Online]. Available: \url{https://github.com/ryrobotics/air_bumper}
\bibitem{whisker} C.~Ye, G.~de~Croon, and S.~Hamaza, ``Whisker-based tactile flight for tiny drones,'' \emph{Nature Communications}, 2025 (arXiv:2510.03119).
\bibitem{lidarsurvey} Y.~Ren, Y.~Cai, H.~Li, N.~Chen, F.~Zhu, L.~Yin, F.~Kong, R.~Li, and F.~Zhang, ``A survey on LiDAR-based autonomous aerial vehicles,'' \emph{IEEE/ASME Trans. Mechatronics}, 2025, doi: 10.1109/TMECH.2025.3599633.
\bibitem{gp} N.~Mohammad, J.~Higgins, and N.~Bezzo, ``A GP-based robust motion planning framework for agile autonomous robot navigation and recovery in unknown environments,'' in \emph{Proc. IEEE Int. Conf. Robot. Autom. (ICRA)}, 2024, pp. 2418--2424, doi: 10.1109/ICRA57147.2024.10610382.
\bibitem{ifar} B.~He, G.~Chen, W.~Wang, J.~Zhang, C.~Fermüller, and Y.~Aloimonos, ``Interactive-FAR: Interactive, fast and adaptable routing for navigation among movable obstacles in complex unknown environments,'' in \emph{Proc. IEEE/RSJ Int. Conf. Intell. Robots Syst. (IROS)}, 2024, pp. 5402--5409, doi: 10.1109/IROS58592.2024.10802169.
\bibitem{ros2} S.~Macenski, T.~Foote, B.~Gerkey, C.~Lalancette, and W.~Woodall, ``Robot Operating System 2: Design, architecture, and uses in the wild,'' \emph{Science Robotics}, vol.~7, no.~66, eabm6074, 2022.
\bibitem{px4} L.~Meier, D.~Honegger, and M.~Pollefeys, ``PX4: A node-based multithreaded open source robotics framework for deeply embedded platforms,'' in \emph{Proc. IEEE Int. Conf. Robot. Autom. (ICRA)}, 2015, pp. 6235--6240.
\bibitem{gazebo} N.~Koenig and A.~Howard, ``Design and use paradigms for Gazebo, an open-source multi-robot simulator,'' in \emph{Proc. IEEE/RSJ Int. Conf. Intell. Robots Syst. (IROS)}, 2004, pp. 2149--2154.
\bibitem{edt} P.~F. Felzenszwalb and D.~P. Huttenlocher, ``Distance transforms of sampled functions,'' \emph{Theory of Computing}, vol.~8, pp. 415--428, 2012.
\bibitem{haddadin} S.~Haddadin, A.~De~Luca, and A.~Albu-Schäffer, ``Robot collisions: A survey on detection, isolation, and identification,'' \emph{IEEE Trans. Robot.}, vol.~33, no.~6, pp. 1292--1312, 2017.
\bibitem{vovk} V.~Vovk, A.~Gammerman, and G.~Shafer, \emph{Algorithmic Learning in a Random World}. New York, NY, USA: Springer, 2005.
\bibitem{angelopoulos} A.~N. Angelopoulos and S.~Bates, ``Conformal prediction: A gentle introduction,'' \emph{Found. Trends Mach. Learn.}, vol.~16, no.~4, pp. 494--591, 2023.
\end{thebibliography}

\end{document}
